\subsection{Ejercicio 5.2: resultados experimentales de YYC y BT}
\label{sec:ejercicio52}

Se procesaron las matrices del punto 5.1 con YYC y BT, primero
con las filas en su orden de generación y después ordenadas de
menor a mayor cantidad de unos. En los empates se conservó el
orden relativo original. El ordenamiento se realizó antes de
iniciar la medición y no forma parte del tiempo del algoritmo.

Las Tablas~\ref{tab:resultadosphi52} y~\ref{tab:resultadosgamma52}
adaptan las Tablas 5 y 6 de \cite{alba2016}: conservan las primeras
cinco columnas y añaden los cuatro tiempos solicitados. Se utiliza
$C=\theta(A,B)$; el caso $N=1$ es la misma matriz en ambas familias
y comparte sus mediciones. Los conteos de testores corresponden
a las propiedades calculadas en 5.1, comprobadas con YYC.

Se registró una ejecución manual por caso, utilizando el tiempo
reportado por el programa, sin traza. En YYC se acumulan los
intervalos de procesamiento de las filas; en BT se mide la
resolución, excluyendo la verificación posterior por minimalidad.
La impresión y exportación quedan fuera de la medición. Los valores
están expresados en milisegundos; los casos sin resultado final
se identifican como \emph{falta ejecutar} y no se completan con
estimaciones. El análisis de estos resultados se presenta en 5.3.

\begin{table}[htbp]
\centering
{\scriptsize
\setlength{\tabcolsep}{3pt}
\begin{tabular}{c c r r r r r r r}
\toprule
\multicolumn{5}{c}{} & \multicolumn{2}{c}{YYC (ms)} & \multicolumn{2}{c}{BT (ms)} \\
\cmidrule(lr){6-7}\cmidrule(lr){8-9}
$N$ & Matriz & Filas & Cols. & $|\Psi^*|$ & Original & Ascendente & Original & Ascendente \\
\midrule
1 & $C$ & 48 & 14 & 14 & 0.2622 & 0.1681 & 0.7632 & 1.027 \\
2 & $\varphi^{2}(C)$ & 48 & 28 & 108 & 13.448 & 5.122 & 1\,169 & 1\,124 \\
3 & $\varphi^{3}(C)$ & 48 & 42 & 396 & 231.593 & 57.589 & 806\,341 & 3\,055\,669 \\
4 & $\varphi^{4}(C)$ & 48 & 56 & 1040 & 1\,847 & 367.384 & \shortstack{falta\\ejecutar} & \shortstack{falta\\ejecutar} \\
5 & $\varphi^{5}(C)$ & 48 & 70 & 2250 & 14\,095 & 2\,192 & \shortstack{falta\\ejecutar} & \shortstack{falta\\ejecutar} \\
\bottomrule
\end{tabular}
}

\caption{Resultados manuales de YYC y BT sobre $\varphi^N(C)$,
en orden original y ascendente. Todos los tiempos están en ms.}
\label{tab:resultadosphi52}
\end{table}

\begin{table}[htbp]
\centering
{\scriptsize
\setlength{\tabcolsep}{3pt}
\begin{tabular}{c c r r r r r r r}
\toprule
\multicolumn{5}{c}{} & \multicolumn{2}{c}{YYC (ms)} & \multicolumn{2}{c}{BT (ms)} \\
\cmidrule(lr){6-7}\cmidrule(lr){8-9}
$N$ & Matriz & Filas & Cols. & $|\Psi^*|$ & Original & Ascendente & Original & Ascendente \\
\midrule
1 & $C$ & 48 & 14 & 14 & 0.2622 & 0.1681 & 0.7632 & 1.027 \\
2 & $\gamma^{2}(C)$ & 96 & 28 & 196 & 34.407 & 11.666 & 31\,476 & 25\,520 \\
3 & $\gamma^{3}(C)$ & 144 & 42 & 2744 & 280.705 & 337.262 & $69\,283\,000^{*}$ & \shortstack{falta\\ejecutar} \\
4 & $\gamma^{4}(C)$ & 192 & 56 & 38416 & 10\,640 & 15\,790 & \shortstack{falta\\ejecutar} & \shortstack{falta\\ejecutar} \\
\bottomrule
\end{tabular}
}

\caption{Resultados manuales de YYC y BT sobre $\gamma^N(C)$,
en orden original y ascendente. Todos los tiempos están en ms.}
\label{tab:resultadosgamma52}
\end{table}

El asterisco de BT original sobre $\gamma^3(C)$ identifica el
registro de 19 h 14 min 43 s, equivalente a 69\,283\,000 ms.
No se verificó la equivalencia de sus condiciones con las de
la serie local; por ello se conserva como registro adicional,
sin utilizarlo para calcular ventajas entre algoritmos.
